% GENERATED FILE. Edit canonical lessons, then run npm run generate:books.
% canonical-source-hash: fnv1a64:b396975eae0524db
% canonical-lessons: ML-C413-suparsa-ceyyuka, ML-C413-abhiprayappetuka, ML-C413-prasamsikkuka, ML-C413-prerippikkuka, ML-C413-prakatippikkuka

\chapter{Opinions and Praise}
\label{ch:ml-opinions-and-praise}
\hlchaptermodality{413}

\begin{chapteropening}
\textbf{By the end of this chapter:} \emph{I can say to recommend, to remark, to praise, to urge and to express.}

Complete the last lesson of chapter 413: I can say to recommend, to remark, to praise, to urge and to express.
\end{chapteropening}

\section[\texorpdfstring{\ml{ശുപാർശ} \ml{ചെയ്യുക} (śupārśa ceyyuka)}{ശുപാർശ ചെയ്യുക (śupārśa ceyyuka)}]{\ml{ശുപാർശ} \ml{ചെയ്യുക} (śupārśa ceyyuka) — to recommend}

\label{lesson:ML-C413-suparsa-ceyyuka}



\begin{quote}
Before the new one: say the Malayalam for generally, then the Malayalam for that means.
\end{quote}

\subsection*{You'll want to know: \ml{ശുപാർശ} \ml{ചെയ്യുക}}
\textbf{\ml{ശുപാർശ} \ml{ചെയ്യുക}} — \emph{śupārśa ceyyuka} — \textquotedblleft{}to recommend\textquotedblright{}.

\begin{quote}
\ml{ഞാൻ} \ml{ആ} \ml{പുസ്തകം} \ml{എല്ലാവർക്കും} \ml{ശുപാർശ} \ml{ചെയ്യുന്നു}.
\end{quote}

\begin{quote}
\emph{ñān ā pustakaṁ ellāvarkkuṁ śupārśa ceyyunnu.}
\end{quote}

\begin{quote}
I recommend that book to everyone.
\end{quote}

Past: \ml{ശുപാർശ} \ml{ചെയ്തു}, \emph{śupārśa ceytu}.

The noun \ml{ശുപാർശ} is a recommendation, also a word put in for someone: \ml{ഒരു} \ml{ശുപാർശക്കത്ത്}, a letter of recommendation.

\begin{cousinweb}
\textbf{\ml{ശുപാർശ} \ml{ചെയ്യുക}} (\emph{śupārśa ceyyuka}) is a phrase verb built on the noun \ml{ശുപാർശ}, \textquotedblleft{}a recommendation\textquotedblright{}, which is usually traced to Persian \emph{sifāriš}, \textquotedblleft{}recommendation\textquotedblright{}; \ml{ചെയ്യുക}, \textquotedblleft{}to do\textquotedblright{}, makes it a verb.
\end{cousinweb}

\begin{grammarlens}[title={What you've built}]
One more word for giving reasons and opinions.
\end{grammarlens}

\subsection*{Your turn}
\begin{itemize}
\raggedright
  \item \emph{Say it:} \emph{śupārśa ceyyuka}
  \item \emph{Say it:} \emph{śupārśa ceyyuka}, once more
  \item \emph{Say it:} \emph{ñān ā pustakaṁ ellāvarkkuṁ śupārśa ceyyunnu.} — I recommend that book to everyone.
  \item \emph{Recall:} say the Malayalam for generally, then the Malayalam for that means, then say \emph{śupārśa ceyyuka} again
\end{itemize}

\begin{tcolorbox}[breakable,colback=teal!4,colframe=teal!35!black,title={Before you move on}]
What does \ml{ശുപാർശ} \ml{ചെയ്യുക} mean? (\textquotedblleft{}To recommend\textquotedblright{}.) Say it once more.
\end{tcolorbox}
\section[\texorpdfstring{\ml{അഭിപ്രായപ്പെടുക} (abhiprāyappeṭuka)}{അഭിപ്രായപ്പെടുക (abhiprāyappeṭuka)}]{\ml{അഭിപ്രായപ്പെടുക} (abhiprāyappeṭuka) — to remark; to opine}

\label{lesson:ML-C413-abhiprayappetuka}



\begin{quote}
Before the new one: say the Malayalam for that means, then the Malayalam for to recommend.
\end{quote}

\subsection*{You'll want to know: \ml{അഭിപ്രായപ്പെടുക}}
\textbf{\ml{അഭിപ്രായപ്പെടുക}} — \emph{abhiprāyappeṭuka} — \textquotedblleft{}to remark; to opine\textquotedblright{}.

\begin{quote}
\ml{ആ} \ml{സിനിമ} \ml{നല്ലതാണെന്ന്} \ml{എല്ലാവരും} \ml{അഭിപ്രായപ്പെട്ടു}.
\end{quote}

\begin{quote}
\emph{ā sinima nallatāṇennŭ ellāvaruṁ abhiprāyappeṭṭu.}
\end{quote}

\begin{quote}
Everyone remarked that the film was good.
\end{quote}

Past: \ml{അഭിപ്രായപ്പെട്ടു}, \emph{abhiprāyappeṭṭu}.

Built on \ml{അഭിപ്രായം}, an opinion. What is said comes before \ml{എന്ന്}.

\begin{cousinweb}
\textbf{\ml{അഭിപ്രായപ്പെടുക}} (\emph{abhiprāyappeṭuka}) is built on \ml{അഭിപ്രായം}, \textquotedblleft{}an opinion\textquotedblright{}, from Sanskrit \emph{abhiprāya}, \textquotedblleft{}intention, opinion\textquotedblright{} (Monier-Williams), with -\ml{പ്പെടുക} from native \ml{പെടുക} (DEDR 3853): to give an opinion.
\end{cousinweb}

\begin{grammarlens}[title={What you've built}]
One more word for giving reasons and opinions.
\end{grammarlens}

\subsection*{Your turn}
\begin{itemize}
\raggedright
  \item \emph{Say it:} \emph{abhiprāyappeṭuka}
  \item \emph{Say it:} \emph{abhiprāyappeṭuka}, once more
  \item \emph{Say it:} \emph{ā sinima nallatāṇennŭ ellāvaruṁ abhiprāyappeṭṭu.} — Everyone remarked that the film was good.
  \item \emph{Recall:} say the Malayalam for that means, then the Malayalam for to recommend, then say \emph{abhiprāyappeṭuka} again
\end{itemize}

\begin{tcolorbox}[breakable,colback=teal!4,colframe=teal!35!black,title={Before you move on}]
What does \ml{അഭിപ്രായപ്പെടുക} mean? (\textquotedblleft{}To remark; to opine\textquotedblright{}.) Say it once more.
\end{tcolorbox}
\section[\texorpdfstring{\ml{പ്രശംസിക്കുക} (praśaṁsikkuka)}{പ്രശംസിക്കുക (praśaṁsikkuka)}]{\ml{പ്രശംസിക്കുക} (praśaṁsikkuka) — to commend; to praise}

\label{lesson:ML-C413-prasamsikkuka}



\begin{quote}
Before the new one: say the Malayalam for to recommend, then the Malayalam for to remark.
\end{quote}

\subsection*{You'll want to know: \ml{പ്രശംസിക്കുക}}
\textbf{\ml{പ്രശംസിക്കുക}} — \emph{praśaṁsikkuka} — \textquotedblleft{}to commend; to praise\textquotedblright{}.

\begin{quote}
\ml{അധ്യാപകൻ} \ml{അവളുടെ} \ml{ഉപന്യാസത്തെ} \ml{പ്രശംസിച്ചു}.
\end{quote}

\begin{quote}
\emph{adhyāpakan avaḷuṭe upanyāsatte praśaṁsiccu.}
\end{quote}

\begin{quote}
The teacher praised her essay.
\end{quote}

Past: \ml{പ്രശംസിച്ചു}, \emph{praśaṁsiccu}.

What is praised takes -\ml{നെ}. The noun is \ml{പ്രശംസ}, praise.

\begin{cousinweb}
\textbf{\ml{പ്രശംസിക്കുക}} (\emph{praśaṁsikkuka}) is built from Sanskrit \emph{praśaṁsā}, \textquotedblleft{}praise, commendation\textquotedblright{} (Monier-Williams), with the Malayalam verb ending -\ml{ിക്കുക}.
\end{cousinweb}

\begin{grammarlens}[title={What you've built}]
One more word for giving reasons and opinions.
\end{grammarlens}

\subsection*{Your turn}
\begin{itemize}
\raggedright
  \item \emph{Say it:} \emph{praśaṁsikkuka}
  \item \emph{Say it:} \emph{praśaṁsikkuka}, once more
  \item \emph{Say it:} \emph{adhyāpakan avaḷuṭe upanyāsatte praśaṁsiccu.} — The teacher praised her essay.
  \item \emph{Recall:} say the Malayalam for to recommend, then the Malayalam for to remark, then say \emph{praśaṁsikkuka} again
\end{itemize}

\begin{tcolorbox}[breakable,colback=teal!4,colframe=teal!35!black,title={Before you move on}]
What does \ml{പ്രശംസിക്കുക} mean? (\textquotedblleft{}To commend; to praise\textquotedblright{}.) Say it once more.
\end{tcolorbox}
\section[\texorpdfstring{\ml{പ്രേരിപ്പിക്കുക} (prērippikkuka)}{പ്രേരിപ്പിക്കുക (prērippikkuka)}]{\ml{പ്രേരിപ്പിക്കുക} (prērippikkuka) — to urge; to motivate}

\label{lesson:ML-C413-prerippikkuka}



\begin{quote}
Before the new one: say the Malayalam for to remark, then the Malayalam for to commend.
\end{quote}

\subsection*{You'll want to know: \ml{പ്രേരിപ്പിക്കുക}}
\textbf{\ml{പ്രേരിപ്പിക്കുക}} — \emph{prērippikkuka} — \textquotedblleft{}to urge; to motivate\textquotedblright{}.

\begin{quote}
\ml{കൂട്ടുകാർ} \ml{അവനെ} \ml{ആ} \ml{ജോലി} \ml{ഉപേക്ഷിക്കാൻ} \ml{പ്രേരിപ്പിച്ചു}.
\end{quote}

\begin{quote}
\emph{kūṭṭukār avane ā jōli upēkṣikkān prērippiccu.}
\end{quote}

\begin{quote}
His friends urged him to give up that job.
\end{quote}

Past: \ml{പ്രേരിപ്പിച്ചു}, \emph{prērippiccu}.

The noun is \ml{പ്രേരണ}, an urging, a motive.

\begin{cousinweb}
\textbf{\ml{പ്രേരിപ്പിക്കുക}} (\emph{prērippikkuka}) is built from Sanskrit \emph{preraṇa}, \textquotedblleft{}driving on, urging\textquotedblright{} (CDIAL 9000), with the Malayalam verb ending -\ml{ിപ്പിക്കുക}. It is the verb of \ml{പ്രേരണ}.
\end{cousinweb}

\begin{grammarlens}[title={What you've built}]
One more word for giving reasons and opinions.
\end{grammarlens}

\subsection*{Your turn}
\begin{itemize}
\raggedright
  \item \emph{Say it:} \emph{prērippikkuka}
  \item \emph{Say it:} \emph{prērippikkuka}, once more
  \item \emph{Say it:} \emph{kūṭṭukār avane ā jōli upēkṣikkān prērippiccu.} — His friends urged him to give up that job.
  \item \emph{Recall:} say the Malayalam for to remark, then the Malayalam for to commend, then say \emph{prērippikkuka} again
\end{itemize}

\begin{tcolorbox}[breakable,colback=teal!4,colframe=teal!35!black,title={Before you move on}]
What does \ml{പ്രേരിപ്പിക്കുക} mean? (\textquotedblleft{}To urge; to motivate\textquotedblright{}.) Say it once more.
\end{tcolorbox}
\section[\texorpdfstring{\ml{പ്രകടിപ്പിക്കുക} (prakaṭippikkuka)}{പ്രകടിപ്പിക്കുക (prakaṭippikkuka)}]{\ml{പ്രകടിപ്പിക്കുക} (prakaṭippikkuka) — to express; to show}

\label{lesson:ML-C413-prakatippikkuka}



\begin{quote}
Before the new one: say the Malayalam for to commend, then the Malayalam for to urge.
\end{quote}

\subsection*{You'll want to know: \ml{പ്രകടിപ്പിക്കുക}}
\textbf{\ml{പ്രകടിപ്പിക്കുക}} — \emph{prakaṭippikkuka} — \textquotedblleft{}to express; to show\textquotedblright{}.

\begin{quote}
\ml{അവൻ} \ml{തന്റെ} \ml{നന്ദി} \ml{ഒരു} \ml{കത്തിലൂടെ} \ml{പ്രകടിപ്പിച്ചു}.
\end{quote}

\begin{quote}
\emph{avan tanṟe nandi oru kattilūṭe prakaṭippiccu.}
\end{quote}

\begin{quote}
He expressed his thanks in a letter.
\end{quote}

Past: \ml{പ്രകടിപ്പിച്ചു}, \emph{prakaṭippiccu}.

Said of feelings and opinions: \ml{അഭിപ്രായം} \ml{പ്രകടിപ്പിക്കുക}, to express an opinion.

\begin{cousinweb}
\textbf{\ml{പ്രകടിപ്പിക്കുക}} (\emph{prakaṭippikkuka}) is built from Sanskrit \emph{prakaṭana}, \textquotedblleft{}bringing to light\textquotedblright{} (CDIAL 8429), with the Malayalam causative verb ending -\ml{ിപ്പിക്കുക}: to bring a feeling to light.
\end{cousinweb}

\begin{grammarlens}[title={What you've built}]
One more word for giving reasons and opinions.
\end{grammarlens}

\subsection*{Your turn}
\begin{itemize}
\raggedright
  \item \emph{Say it:} \emph{prakaṭippikkuka}
  \item \emph{Say it:} \emph{prakaṭippikkuka}, once more
  \item \emph{Say it:} \emph{avan tanṟe nandi oru kattilūṭe prakaṭippiccu.} — He expressed his thanks in a letter.
  \item \emph{Recall:} say the Malayalam for to commend, then the Malayalam for to urge, then say \emph{prakaṭippikkuka} again
\end{itemize}

\begin{tcolorbox}[breakable,colback=teal!4,colframe=teal!35!black,title={Before you move on}]
What does \ml{പ്രകടിപ്പിക്കുക} mean? (\textquotedblleft{}To express; to show\textquotedblright{}.) Say it once more.
\end{tcolorbox}
